% ============================================================================
% ABSTRAK: kalimat hasil dan implikasi menunggu eksperimen E1--E4.
% ============================================================================
\begin{abstract}
Perhitungan PPh 21 harus mengikuti perubahan regulasi dan kebijakan perusahaan, tetapi aturannya umumnya tertanam di kode program atau \en{spreadsheet}, sehingga setiap perubahan menuntut pemrograman ulang dan dapat menggeser hasil periode lama tanpa terdeteksi. Penelitian formalisasi hukum pajak berbantuan \en{large language model} (LLM) membangun basis aturan dari nol untuk satu keadaan hukum, sehingga belum diketahui apakah pembaruan berulang oleh LLM menjaga hasil periode lama dan apakah kesalahannya merambat dari satu perubahan ke perubahan berikutnya. Penelitian ini mengusulkan \en{knowledge-based system} adaptif yang menyimpan pengetahuan sebagai \en{legal knowledge graph} temporal dua lapis, yaitu lapisan regulasi dan lapisan kebijakan perusahaan, dan menjalankan seluruh perhitungannya dengan \en{inference engine} deterministik. Pembaruan \en{knowledge base} disusun LLM dari dokumen regulasi atau kebijakan, lalu disaring lima gerbang verifikasi dan dipilih lewat konsensus perilaku. Evaluasinya memutar ulang perubahan regulasi PPh 21 pegawai tetap yang nyata pada 2023--2026, termasuk pergantian metode ke tarif efektif rata-rata, menguji perubahan yang belum pernah dilihat perancang, dan membandingkan usulan dengan pembanding serta ablasi. \isian{Kalimat hasil: RA, SER dan CAR dengan selang kepercayaan 95\%, OAR, dan \en{recall} analisis dampak.} \isian{Kalimat implikasi.}
\end{abstract}

\begin{IEEEkeywords}
Akuisisi pengetahuan, \en{knowledge-based system}, \en{knowledge graph}, \en{large language model}, perubahan regulasi, PPh Pasal 21, \en{rules as code}
\end{IEEEkeywords}
